\subsection{Small support lattices and complement failures}

\begin{proposition}[A ten-element ortholattice]\label{prop:ten-lattice}
Let $I=\{i,i',j,k\}$ and
\[
 T=\{A\subseteq I:(i\in A\Longleftrightarrow i'\in A)\}
       \cup\{\{i\},\{i',j,k\}\}.
\]
Inclusion and set complement make $T$ a ten-element ortholattice
which is not orthomodular. With labels
\[
\begin{array}{c|ccccc}
&0&a'&a&b&c\\\hline
&\varnothing&\{i\}&\{i,i'\}&\{j\}&\{k\}\\[3pt]
&A&A'&B&C&1\\\hline
&\{j,k\}&\{i',j,k\}&\{i,i',k\}&\{i,i',j\}&I,
\end{array}
\]
its covers are
\[
\begin{gathered}
0\prec a'\prec a,\quad 0\prec b,c,\quad a\prec B,C,\\
b\prec A,C,\quad c\prec A,B,\quad A\prec A'\prec1,
\quad B,C\prec1.
\end{gathered}
\]
In particular $c\not\leq C$.
\end{proposition}
\begin{proof}
The paired subsets form an eight-element Boolean cube; the other
two subsets are new. If an intersection of two members belongs
to $T$, it is their meet. Otherwise it contains $i'$ but not $i$
and is a proper subset of $\{i',j,k\}$. Deleting $i'$ gives its
largest allowed subset: no allowed subset of that intersection
can contain $i'$. Thus every meet exists. Complement preserves
$T$ and reverses inclusion, so it gives every join by duality.
Complementary members have meet $0$ and join $1$.

For $x=\{i\}<y=\{i,i'\}$, the intersection with $x^\perp$
is $\{i'\}$ and its allowed meet is empty. Hence
$x\vee(y\wedge x^\perp)=x\ne y$.
The minimal nonzero members are $a',b,c$; the only immediate
successor of $a'$ is $a$. Listing the minimal allowed enlargements
of $a,b,c,A$ gives the displayed covers, and complementation
gives the upper covers. Finally $k\notin\{i,i',j\}$.
\end{proof}

This is the ordinary-powerset reconstruction of the formula in
\cite[part~9, archive sheet~112, PDF p.~113]{Groth}. The handwritten
powerset mark and overwritten drawing do not justify treating every
visible stroke as a covering relation. Reading the mark instead as
``three-element subsets'' yields just three incomparable sets,
not this ortholattice; the mathematical convention is explicit here.

\begin{proposition}[Identity-incidence graph failures]\label{prop:graph-failures}
Give a simple graph identity incidence and internal refinement,
with adjacency as disjunction and all cliques as figures. This
is a strong system with nonempty open-lieu sets and local detection.
For the path $a'\! -\! a\! -\! b\! -\! b'$, its support lattice is
$O_6$, and $F=\{a,b\}$ satisfies
\[
 \cl F\cap N(a)=\{a',b\}\ne\{b\}=\cl(F\setminus\{a\}).
\]
For the five-cycle, singleton supports are singletons, yet the
same relative-complement identity fails. Thus it can fail even
when spatial refinement is reconstructed from singleton supports.
\end{proposition}
\begin{proof}
Principal deployments are singletons, so the carrier and exchange
conditions are identities, and principal disjunction and heredity
are automatic. Every object is a lieu, with $U_x=\{x\}$; all local
tests are exactly the graph tests. Figure refinement is inclusion
and a subdivision is equality: a missing vertex is disjoint from
the smaller clique but not from the larger one containing it.

The path neighborhoods are
\[
 N(a')=\{a\},\quad N(a)=\{a',b\},\quad
 N(b)=\{a,b'\},\quad N(b')=\{b\}.
\]
All intersections of these neighborhoods give exactly the six
closed sets
\[
 \varnothing,\quad E,\quad\{a\},\quad\{b\},\quad\{a',b\},\quad\{a,b'\},
\]
whose order and complements are $O_6$. The adjacent pair has
empty polar and hence full closure, proving its displayed failure.

Number the five-cycle modulo five. The common neighbors of
$i-1$ and $i+1$ consist of $i$, so $\cl\{i\}=\{i\}$.
For $F=\{0,1\}$, its polar is empty and
\[
 \cl F\cap N(0)=\{1,4\}\ne\{1\}=\cl\{1\}.
\]
Inclusion carriers between singleton labelled deployments exist
exactly for identical vertices, so the spatial reconstruction
is identity as asserted. In this edge figure the meet of the
two singleton supports is still zero; the complement failure
alone does not establish a meet failure.
\end{proof}

These identity systems also satisfy the following explicit
witness-absorption property: if $G\ll F$ and a figure is disjoint
from $G$ but not from $F$, some figure refining $F$ has that
property. Indeed $G\subsetneq F$, and a missing vertex of the
clique $F$ is such a witness. For the five-cycle the open support
of each stratum is its singleton, so it is absorbed by that
stratum's refinement ideal as well.

\subsection{External support functions}
For this paragraph an external support assignment on a set $P$
is \emph{normalized} if its supports cover $P$, and \emph{faithful}
if their disjointness detects primitive disjunction. For identity
incidence the open and full external supports coincide. The
stronger \emph{full faithfulness} condition in
\cite[archive sheets~84--85, PDF pp.~85--86]{Groth8} requires
\[
 \bigcup_{z\perp F}|z|^\circ=P\setminus|F|
 \quad\text{for every figure }F.
\]
Here $|F|$ is the union of its stratum supports.

\begin{proposition}[Faithful does not imply fully faithful]
\label{prop:external-faithful}
On two strata $a,b$, take identity incidence, $a\rho b$ besides
identities, and empty disjunction. With $P=\{p,q\}$, the assignment
$|a|=\{p\}$ and $|b|=P$ is normalized, monotone under refinement,
faithful, and separates the points of $P$, but is not fully faithful.
\end{proposition}
\begin{proof}
All deployments are singletons; the proper internal relation
is a partial-order relation and has no nontrivial exchange case.
Thus the data form a strong system, with figures
$\varnothing,\{a\},\{b\}$. All external supports are nonempty and
intersect each other, as required by empty disjunction. They cover
$P$; membership in $|a|$ distinguishes $p$ and $q$.
For $F=\{a\}$ there is no disjoint stratum, so the left side of
full faithfulness is empty and the right side is $\{q\}$.
Notice that inclusion of supports does reflect refinement in
this example; that is a different condition.
\end{proof}

\subsection{Completing labelled spatial deployments}
We prove Theorem~\ref{thm:spatial-completion}. A carrier, when it
exists, is unique: a nonempty label cannot lie inside two distinct
orthogonal labels in the target deployment. Identities and
composites are inclusion carriers. If carriers exist in both
directions, their composite sends each label $S$ to a label of
the same orthogonal family containing $S$. That label must be $S$.
The intermediate inclusion is therefore also equality. Both
carriers fix the actual labels and give identical ordered
labelled deployments, so their source strata coincide by the
injectivity hypothesis. This proves antisymmetry, including for
the relation allowing carriers that are not internal.

For a carrier $h:D_x\to D_y$, let $S_z=h(S_x)$. Order preservation
puts every image below $S_z$. That lower ideal of $D_y$ is precisely
$D_z$, and hence $h$ factors as an internal carrier to $D_z$ followed
by its incidence inclusion. Uniqueness of carriers gives uniqueness
of this factor. Restricting to the principal ideal of a face of
$x$ gives the exchange factor, and all composition compatibilities
follow from uniqueness. The internal relation is reflexive and
transitive, and the preceding antisymmetry argument makes it a
partial order. Pairwise orthogonality in each $D_x$ is precisely
principal disjunction. If $x'\rho x$, then $S_{x'}\subseteq S_x$,
so orthogonality to another label is inherited by $S_{x'}$.
All strong factorization conditions are proved.

For the weaker reading, use the five-cycle with incidence
$0,2<1$ and identities otherwise. Its singleton closed supports
are singleton sets. Every principal labelled deployment is a
family of distinct, disjoint nonempty sets, and all ordered
labelled deployments are different. Both proper incidences
join orthogonal vertices. Nevertheless $0$ and $2$ are not
orthogonal, although both lie in the principal ideal of $1$.
Thus no choice of internal refinement can make these data a
strong system: its principal-disjunction condition already fails.
The orthogonal spatial target in the theorem is essential.

\subsection{An eight-stratum meet failure}
The support calculus tested here is the one formed from open
stratum supports in \cite[archive sheets~92--93, PDF pp.~93--94]{Groth7};
the corresponding absorption condition appears on archive
sheet~86, PDF p.~87. We state the mathematical hypotheses explicitly.
We distinguish primitive disjunction $\delta$ from disjointness
of whole principal figures. Put $x\mathrel{\tau}y$ when $[x]$
and $[y]$ are disjoint figures whose union is a figure. Write
$N_\tau$ and $\cl_\tau$ for its polarity and closure. Define the
open support of a stratum by
\[
 U_x=\cl_\tau([x])\cap N_\tau([x]\setminus\{x\}).
\]
For a fixed figure $F$ and any index subset $A\subseteq F$, put
$S_A=\cl_\tau(\bigcup_{x\in A}U_x)$. The index subset need not
itself be incidence-closed.

\begin{proposition}[Failure within one figure]\label{prop:eight-meet}
Let $M=\{a,b,X,p,e,w,h,g\}$, with incidence $a,b<X$, internal
refinement $p\rho X$, and no other proper generating relations.
Give the seven point strata the disjunction edges
\[
 ab,\ ap,\ bp,\ ae,\ bw,\ pw,\ eh,\ hg,\ gw,
\]
and add $Xa,Xb,Xw$ to primitive disjunction. Use all
incidence-closed primitive cliques as figures. This is a strong
system with nonempty open-lieu sets, local detection, and the
witness-absorption property stated above. Every open support
$U_x$ lies in the refinement ideal of $x$. Nevertheless, in
$F=\{a,b,X\}$,
\[
\begin{aligned}
 S_F\cap N_\tau(S_{\{a\}})&=\{b,p,e\}
       \ne\{b,p\}=S_{\{b,X\}},\\
 S_{\{a,X\}}\cap S_{\{a,b\}}&=\{a,w\}
       \ne\{a\}=S_{\{a\}}.
\end{aligned}
\]
Thus both relative complementation and nonempty-meet preservation
fail for the support calculus of the strata of this one figure.
\end{proposition}
\begin{proof}
The only nonsingleton principal deployment is $\{a,b,X\}$, and
its members are pairwise primitively disjoint. The proper internal
arrow starts at a singleton deployment, so exchange is immediate.
Its heredity requirement holds since $p$ is disjoint from $a,b,w$.
These observations verify the strong conditions. The lieux are
the seven point strata; $X$ has open-lieu set $\{p\}$, while
each point has its singleton. Primitive disjunction of $X$ agrees
with that of $p$, so local detection holds in every case.

For witness absorption, a figure not containing $X$ is a point
clique and every refining figure is its subset. Any missing point
is a witness. A figure containing $X$ must be $F$: no stratum is
disjoint from all of $a,b,X$. A proper figure refining $F$ is
therefore a subset of the point triangle $\{a,b,p\}$. If $a$ or
$b$ is missing, use that point as witness. If the refining figure
is $\{a,b\}$, use $\{p\}$. The remaining case is $\{a,b,p\}$;
no nonempty figure is disjoint from it, since the three point
neighborhoods have empty intersection. The separation premise
then fails. The case of the whole $F$ has a false premise too.

The whole-figure relation $\tau$ has the nine displayed point
edges and makes $X$ isolated. Its point neighborhoods are
\[
\begin{array}{c|cccc}
x&a&b&p&e\\\hline
N_\tau(x)&\{b,p,e\}&\{a,p,w\}&\{a,b,w\}&\{a,h\}\\[3pt]
x&w&h&g&\\\hline
N_\tau(x)&\{b,p,g\}&\{e,g\}&\{w,h\}&
\end{array}
\]
Intersecting neighborhoods shows that every point has singleton
closed support. Since $X$ is isolated, $\cl_\tau([X])=M$ and
$N_\tau(\{a,b\})=\{p\}$. Thus $U_X=\{p\}$, and every other
$U_x=\{x\}$; these are all absorbed by their refinement ideals.
Finally direct intersections of this table give
\[
\begin{gathered}
 S_F=M,\quad S_{\{a\}}=\{a\},\quad
 S_{\{b,X\}}=\{b,p\},\\
 S_{\{a,X\}}=\{a,p,w\},\qquad
 S_{\{a,b\}}=\{a,b,w\}.
\end{gathered}
\]
The two displayed failures follow. This calculation uses $\tau$
and the open supports, throughout; substituting primitive
$\delta$ changes some closed sets.
\end{proof}
